\section{A compact transition and its finite lifts}\label{sec:geometry}

We use the four-node transition denoted $X_9$ in \cite{WuebbenMirror}.
For clarity, the numerical label only identifies the toric example
specified below. Its source fan, divisor basis and smoothing fan are
printed in Appendix~\ref{app:data}. The general anticanonical
hypersurfaces are taken nondegenerate with respect to their Newton
polytopes, with nonzero vertex coefficients.

\subsection{The geometric hypotheses and class marking}
The nodal model $X'_9$ is projective, normal, has trivial dualizing
sheaf and $H^1(\mathcal O)=H^2(\mathcal O)=0$. Its four nodes admit
a projective small resolution $Y_-$. We use the nodal model and the
integral divisor-evaluation injection of
\cite[Proposition 3.5, Lemma 9.7 and Section 11.1]{WuebbenMirror}.
The source $Y_+$ used for the explicit calculation differs from
$Y_-$ by one flop: the original square diagonals $23,24$ are replaced
by $67,24$. The relevant square circuits are
$v_2+v_3=v_6+v_7$ and $v_2+v_4=v_5+v_7$.
The pentagon with vertices $v_3,v_4,v_5,v_6,v_8$ is subdivided
at its interior ray $v_9$ in both resolutions. Our ray indexing
has the pentagon center at $9$ and the divisor missing the
hypersurface at $10$; these two indices are interchanged in
\cite{WuebbenMirror}.

Let $H_1,\ldots,H_6$ be the integral nef basis on $Y_+$ given in
Appendix~\ref{app:data}, and let $b_1,\ldots,b_6$ be the dual curve
basis. Identify the full integral lattices of $Y_-$ and $Y_+$ by the
common blow-up, as at the end of Section~\ref{sec:specialization}.
The four exceptional classes in these coordinates are
\[
\begin{array}{c|rrrr}
 &q_1&q_2&q_3&q_4\\ \hline
Y_-&b_6&-b_5&b_1+b_5+b_6&b_1\\
Y_+&b_6& b_5&b_1+b_5+b_6&b_1.
\end{array}
\]
Their relation on $Y_-$ is $(1,-1,-1,1)$. More intrinsically, the
four divisor-evaluation columns on $Y_-$ are
\[
 e_6+e_8-e_3-e_9,\quad e_6+e_7-e_2-e_3,\quad
 e_5+e_8-e_4-e_9,\quad e_5+e_7-e_2-e_4.
\]
On the first three columns, rows $3,4,9$ have determinant $-1$.
Thus their span is primitive in the ambient integer lattice and in
the full curve lattice, which injects into it. This proves that the
source and its exceptional quotient are torsion-free; it is more
than an ambient rational rank calculation.

The nowhere-zero exceptional relation supplies a global first-order
smoothing nonzero at every node. The abstract deformation space is
smooth by the nodal case of Friedman's unobstructedness theorem
\cite[Theorem 1.6]{Friedman}: the singularities are rational,
Du~Bois and local complete intersections, and the projective
resolution satisfies the $\partial\bar\partial$-lemma.
To obtain an algebraic family of the kind required in
Theorem~\ref{thm:specialization}, choose a sufficiently high ample
embedding of $X'_9$. The vanishing of $H^2(\mathcal O)$ lifts its
line bundle, and the vanishing of $H^1(L^m)$ for large $m$ lifts a
projective frame. Thus embedded deformations map smoothly to
abstract deformations. Equivalently one can use the Hilbert-scheme
construction in \cite[Theorem 4.1]{FMR}. Its local sheaf-generation
condition follows from a tangent nonzero at each node; surjectivity
onto all local smoothing parameters simultaneously is not needed.
Choose a smooth algebraic curve in this Hilbert scheme with the
specified tangent. The restricted universal family is projective,
and at each node has equation
$\sum_{j=1}^4 z_j^2=s_i(t)$ with $s_i'(0)\ne0$.
Its total fourfold is smooth. Relative duality and base change extend
the nowhere-vanishing canonical section near the central fiber, so
the nearby fibers have trivial canonical bundle.

The toric smoothing model in Appendix~\ref{app:data} belongs to this
smooth deformation component, by \cite[Theorem 9.25]{WuebbenMirror}.
Fix transport along a path in its smooth-fiber locus, and choose the
marking for which the smoothing divisors $K_1,K_2,K_3$ descend from
$H_2,H_3,H_4$. Then the full integral class map is
\begin{equation}\label{eq:pi}
 \pi(n_1,\ldots,n_6)=(n_2,n_3,n_4).
\end{equation}
The exceptional subgroup is $\ZZ b_1\oplus\ZZ b_5\oplus\ZZ b_6$.
These class identifications retain the chosen transport path.

\subsection{An explicit bound on the lifts}
Blow up the ambient pair cones $24,67,49,39$, appending the sums of
their rays as rays $11,12,13,14$. These four ambient centers meet
the hypersurface transversely in its four disjoint exceptional
curves. Their possible mutual intersections are outside the
hypersurface. The resulting smooth 47-cone fan therefore contains
$W=\Bl_{\coprod C_i}Y_+$ as the strict transform. The pulled-back
original anticanonical equation defines this strict transform, since
none of the ambient centers is contained in the hypersurface. Its
divisor coefficients are one on the eleven original rays and two
on the four new rays; its Newton polytope is unchanged.
Intersecting its ambient Mori cone with zero degree on the four
exceptional divisors and on $D_{10}$ gives the necessary inequalities
\begin{equation}\label{eq:outer}
 \begin{gathered}
 n_a\ge0,\quad n_6\ge n_1,n_5,\quad
 n_1\le n_2+n_3,\quad n_1\le n_2+n_5,\\
 n_6\le n_2+n_3+n_5,\qquad
 n_1+n_6\le2n_2+n_3+n_5.
 \end{gathered}
\end{equation}
These follow from the 94 invariant walls of the subdivided fan.

One further inequality is needed. On the common blow-up the divisor
\[
 \mathcal D=\phi^*(2H_2+3H_3+H_4-H_6)-\sum_i E_i
\]
is nef. To check this, its coefficients on rays $0,\ldots,14$ are
\[
 d=(0,0,2,0,0,7,7,-5,-5,3,0,1,1,2,2).
\]
On the single open toric neighborhood $\widetilde P\setminus D_{10}$
of $W$, the characters in Table~\ref{tab:characters} give global
sections of $2\mathcal D$. For each cone $\sigma$ maximal among those
whose orbits can meet $W$, the displayed character $m$ satisfies
\[
 \langle m,v_i\rangle+2d_i\ge0\quad(i\ne10),\qquad
 \langle m,v_i\rangle+2d_i=0\quad(i\in\sigma).
\]
Here the Newton faces are those of the pulled-back original equation.
There are 99 cones with nonsingleton faces, each a face of one of the
41 listed cones. The remaining orbits have singleton Newton faces
and miss the hypersurface. The sections are global on the same open
set and generate along $W$; hence $2\mathcal D|_W$ is globally
generated. Pairing with $j\gamma$ gives
\begin{equation}\label{eq:last-bound}
 n_6\le2n_2+3n_3+n_4.
\end{equation}

Let $\beta=(1,1,0)$ on the smoothing. By
Theorem~\ref{thm:specialization}, all possible nonzero normalized
source coefficients over $d\beta$ have class $(a,d,d,0,b,c)$, where
\begin{equation}\label{eq:candidates}
 0\le a\le2d,\qquad b\ge\max(0,a-d),\qquad
 \max(a,b)\le c\le\min(5d,2d+b,3d+b-a).
\end{equation}
There are 37 integer solutions for $d=1$, and 183 for $d=2$.
For $d=1$ they are listed compactly by the permitted values of $c$:
\[
\begin{array}{c|rrrrrr}
a\backslash b&0&1&2&3&4&5\\ \hline
0&0\!:\!2&1\!:\!3&2\!:\!4&3\!:\!5&4\!:\!5&5\\
1&1\!:\!2&1\!:\!3&2\!:\!4&3\!:\!5&4\!:\!5&5\\
2&-&2&2\!:\!3&3\!:\!4&4\!:\!5&5.
\end{array}
\]
These lists are necessary-condition supersets; effectivity and
nonvanishing are not asserted for each entry. The same inequalities
give nine candidates in bidegree $(n_2,n_3)=(1,0)$ and ten in
bidegree $(0,1)$, with $n_4=0$.

\subsection{The plane on the smoothing}
\begin{proposition}\label{prop:plane}
The smoothing $S$ contains a plane $D\simeq\PP^2$ with
$[D]=-2K_1-K_2+K_3$ and normal bundle $\mathcal O(-3)$.
Its line class is primitive and equals $\beta=(1,1,0)$.
Every effective cycle of class $d\beta$, $d>0$, is supported on $D$.
\end{proposition}
\begin{proof}
Use the smoothing rays $r_i$ of Appendix~\ref{app:data} and write
$T_i$ for their divisors. The open set away from $T_4\cup T_5$ has
cones $1236,0236,0123$, all unimodular, and ray relation
$r_0+r_1-3r_2+2r_3+r_6=0$. It is
$\Tot_{\PP^2}(\mathcal O(-3)\oplus\mathcal O(2))$, with fiber
coordinates $u,v$, and $T_2=\{u=0\}$. The anticanonical equation
on $u=0$ is $av+q_2(x_0,x_1,x_6)=0$. Its Newton face consists of
the six quadratics and the vertex monomial $v$, so $a\ne0$.
There is no cone containing $24$; the only intersection of $T_2$
with $T_5$ is the fixed point of the cone $01256$, missed by the
hypersurface. Thus all of $D=S\cap T_2$ lies in this open set and
is the graph of $v=-q_2/a$ over the plane.

The charge matrix gives $T_0=K_1$, $T_4=K_3$, $T_5=K_2-K_1$,
and $T_2=-2K_1-K_2+K_3$. Hence $K_1|_D=K_2|_D=h$ and
$K_3|_D=0$, proving the class and normal bundle statements.
Subdivide $01256$ by either $\{1256,0256,0125\}$ or
$\{0156,0126\}$. Both modifications are isomorphisms near $S$.
Convex support functions give $T_0,T_4$ nef on the first and
$T_5,T_4$ nef on the second. Thus every effective class
$(\ell,m,n)$ on $S$ satisfies $\ell\ge0$, $m-\ell\ge0$, $n\ge0$.
Every component of a cycle in $d\beta$ has class $(a,a,0)$,
$a>0$, and intersects $D$ in $-3a$. An effective Cartier divisor
has nonnegative degree on a curve outside it, so all components
lie in $D$. Finally $K_1\cdot\beta=1$ proves primitivity.
\end{proof}
